% GENERATED FILE - do not hand-edit.
% Source of truth: data/calculus-reviewed/c-6.json in the zenith-website repo.
% Regenerate with: python3 scripts/build_calculus_tex.py
\documentclass[11pt]{article}
\usepackage[utf8]{inputenc}
\usepackage[T1]{fontenc}
\usepackage{amsmath,amssymb}
\usepackage{graphicx}
\usepackage[margin=1in]{geometry}
\graphicspath{{../..//}}
\setlength{\parindent}{0pt}
\setlength{\parskip}{0.6em}
\begin{document}
\section*{Chapter 6: Definite Integrals --- Practice Exercises C6}

\bigskip\noindent\textbf{A1.}\ $\int_{-1}^1(x^2-x-1)\,dx=$
\begin{list}{}{\setlength{\leftmargin}{3em}\setlength{\labelwidth}{2.5em}\setlength{\itemsep}{0pt}\setlength{\topsep}{2pt}}
\item[(A)] $2/3$
\item[(B)] $0$
\item[(C)] $-4/3$
\item[(D)] $-2$
\end{list}

\bigskip\noindent\textbf{A2.}\ $\int_1^2\frac{3x-1}{3x}\,dx=$
\begin{list}{}{\setlength{\leftmargin}{3em}\setlength{\labelwidth}{2.5em}\setlength{\itemsep}{0pt}\setlength{\topsep}{2pt}}
\item[(A)] $3/4$
\item[(B)] $1-\frac13\ln2$
\item[(C)] $1-\ln2$
\item[(D)] $-\frac13\ln2$
\end{list}

\bigskip\noindent\textbf{A3.}\ $\int_0^3\frac{dt}{\sqrt{4-t}}=$
\begin{list}{}{\setlength{\leftmargin}{3em}\setlength{\labelwidth}{2.5em}\setlength{\itemsep}{0pt}\setlength{\topsep}{2pt}}
\item[(A)] $-2$
\item[(B)] $4$
\item[(C)] $-1$
\item[(D)] $2$
\end{list}

\bigskip\noindent\textbf{A4.}\ $\int_{-1}^0\sqrt{3u+4}\,du=$
\begin{list}{}{\setlength{\leftmargin}{3em}\setlength{\labelwidth}{2.5em}\setlength{\itemsep}{0pt}\setlength{\topsep}{2pt}}
\item[(A)] $2$
\item[(B)] $14/9$
\item[(C)] $14/3$
\item[(D)] $7/2$
\end{list}

\bigskip\noindent\textbf{A5.}\ $\int_2^3\frac{dy}{2y-3}=$
\begin{list}{}{\setlength{\leftmargin}{3em}\setlength{\labelwidth}{2.5em}\setlength{\itemsep}{0pt}\setlength{\topsep}{2pt}}
\item[(A)] $\ln3$
\item[(B)] $\frac12\ln(3/2)$
\item[(C)] $16/9$
\item[(D)] $\ln\sqrt3$
\end{list}

\bigskip\noindent\textbf{A6.}\ $\int_0^{\sqrt3}\frac{x}{\sqrt{4-x^2}}\,dx=$
\begin{list}{}{\setlength{\leftmargin}{3em}\setlength{\labelwidth}{2.5em}\setlength{\itemsep}{0pt}\setlength{\topsep}{2pt}}
\item[(A)] $1$
\item[(B)] $\pi/6$
\item[(C)] $-1$
\item[(D)] $2$
\end{list}

\bigskip\noindent\textbf{A7.}\ $\int_0^1(2t-1)^3\,dt=$
\begin{list}{}{\setlength{\leftmargin}{3em}\setlength{\labelwidth}{2.5em}\setlength{\itemsep}{0pt}\setlength{\topsep}{2pt}}
\item[(A)] $1/4$
\item[(B)] $1/2$
\item[(C)] $0$
\item[(D)] $4$
\end{list}

\bigskip\noindent\textbf{A8.}\ $\int_4^9\frac{2+x}{2\sqrt{x}}\,dx=$
\begin{list}{}{\setlength{\leftmargin}{3em}\setlength{\labelwidth}{2.5em}\setlength{\itemsep}{0pt}\setlength{\topsep}{2pt}}
\item[(A)] $25/3$
\item[(B)] $41/3$
\item[(C)] $100/3$
\item[(D)] $5/3$
\end{list}

\bigskip\noindent\textbf{A9.}\ $\int_{-3}^3\frac{dx}{9+x^2}=$
\begin{list}{}{\setlength{\leftmargin}{3em}\setlength{\labelwidth}{2.5em}\setlength{\itemsep}{0pt}\setlength{\topsep}{2pt}}
\item[(A)] $\pi/2$
\item[(B)] $0$
\item[(C)] $\pi/6$
\item[(D)] $-\pi/2$
\end{list}

\bigskip\noindent\textbf{A10.}\ $\int_0^1e^{-x}\,dx=$
\begin{list}{}{\setlength{\leftmargin}{3em}\setlength{\labelwidth}{2.5em}\setlength{\itemsep}{0pt}\setlength{\topsep}{2pt}}
\item[(A)] $1/e-1$
\item[(B)] $-1/e$
\item[(C)] $1-1/e$
\item[(D)] $1/e$
\end{list}

\bigskip\noindent\textbf{A11.}\ $\int_0^1xe^{x^2}\,dx=$
\begin{list}{}{\setlength{\leftmargin}{3em}\setlength{\labelwidth}{2.5em}\setlength{\itemsep}{0pt}\setlength{\topsep}{2pt}}
\item[(A)] $e-1$
\item[(B)] $(e-1)/2$
\item[(C)] $2(e-1)$
\item[(D)] $e/2-1$
\end{list}

\bigskip\noindent\textbf{A12.}\ $\int_0^{\pi/4}\sin(2\theta)\,d\theta=$
\begin{list}{}{\setlength{\leftmargin}{3em}\setlength{\labelwidth}{2.5em}\setlength{\itemsep}{0pt}\setlength{\topsep}{2pt}}
\item[(A)] $2$
\item[(B)] $1/2$
\item[(C)] $-1$
\item[(D)] $-1/2$
\end{list}

\bigskip\noindent\textbf{A13.}\ $\int_1^2\frac{dz}{3-z}=$
\begin{list}{}{\setlength{\leftmargin}{3em}\setlength{\labelwidth}{2.5em}\setlength{\itemsep}{0pt}\setlength{\topsep}{2pt}}
\item[(A)] $-\ln2$
\item[(B)] $3/4$
\item[(C)] $\frac12\ln2$
\item[(D)] $\ln2$
\end{list}

\bigskip\noindent\textbf{A14.}\ If $x=2\sin\theta$, then $\int_1^2\frac{\sqrt{4-x^2}}x\,dx$ is equivalent to
\begin{list}{}{\setlength{\leftmargin}{3em}\setlength{\labelwidth}{2.5em}\setlength{\itemsep}{0pt}\setlength{\topsep}{2pt}}
\item[(A)] $2\int_1^2\frac{\cos^2\theta}{\sin\theta}\,d\theta$
\item[(B)] $\int_{\pi/6}^{\pi/2}\frac{\cos\theta}{\sin\theta}\,d\theta$
\item[(C)] $2\int_{\pi/6}^{\pi/2}\frac{\cos^2\theta}{\sin\theta}\,d\theta$
\item[(D)] $\int_1^2\frac{\cos\theta}{\sin\theta}\,d\theta$
\end{list}

\bigskip\noindent\textbf{A15.}\ $\int_0^\pi\cos^2\theta\sin\theta\,d\theta=$
\begin{list}{}{\setlength{\leftmargin}{3em}\setlength{\labelwidth}{2.5em}\setlength{\itemsep}{0pt}\setlength{\topsep}{2pt}}
\item[(A)] $-2/3$
\item[(B)] $1$
\item[(C)] $2/3$
\item[(D)] $0$
\end{list}

\bigskip\noindent\textbf{A16.}\ $\int_1^e\frac{\ln x}{x}\,dx=$
\begin{list}{}{\setlength{\leftmargin}{3em}\setlength{\labelwidth}{2.5em}\setlength{\itemsep}{0pt}\setlength{\topsep}{2pt}}
\item[(A)] $1/2$
\item[(B)] $(e^2-1)/2$
\item[(C)] $1$
\item[(D)] $e-1$
\end{list}

\bigskip\noindent\textbf{A17.}\ $\int_0^1xe^x\,dx=$
\begin{list}{}{\setlength{\leftmargin}{3em}\setlength{\labelwidth}{2.5em}\setlength{\itemsep}{0pt}\setlength{\topsep}{2pt}}
\item[(A)] $-1$
\item[(B)] $e+1$
\item[(C)] $1$
\item[(D)] $e-1$
\end{list}

\bigskip\noindent\textbf{A18.}\ $\int_0^{\pi/6}\frac{\cos\theta}{1+2\sin\theta}\,d\theta=$
\begin{list}{}{\setlength{\leftmargin}{3em}\setlength{\labelwidth}{2.5em}\setlength{\itemsep}{0pt}\setlength{\topsep}{2pt}}
\item[(A)] $\ln2$
\item[(B)] $3/8$
\item[(C)] $-\frac12\ln2$
\item[(D)] $\ln\sqrt2$
\end{list}

\bigskip\noindent\textbf{A19.}\ $\int_{\sqrt2}^2\frac{u}{u^2-1}\,du=$
\begin{list}{}{\setlength{\leftmargin}{3em}\setlength{\labelwidth}{2.5em}\setlength{\itemsep}{0pt}\setlength{\topsep}{2pt}}
\item[(A)] $\ln\sqrt3$
\item[(B)] $8/9$
\item[(C)] $\ln(3/2)$
\item[(D)] $\ln3$
\end{list}

\bigskip\noindent\textbf{A20.}\ $\int_{\sqrt2}^2\frac{u}{(u^2-1)^2}\,du=$
\begin{list}{}{\setlength{\leftmargin}{3em}\setlength{\labelwidth}{2.5em}\setlength{\itemsep}{0pt}\setlength{\topsep}{2pt}}
\item[(A)] $-1/3$
\item[(B)] $-2/3$
\item[(C)] $2/3$
\item[(D)] $1/3$
\end{list}

\bigskip\noindent\textbf{A21.}\ $\int_{\pi/12}^{\pi/4}\frac{\cos2x}{\sin^22x}\,dx=$
\begin{list}{}{\setlength{\leftmargin}{3em}\setlength{\labelwidth}{2.5em}\setlength{\itemsep}{0pt}\setlength{\topsep}{2pt}}
\item[(A)] $1$
\item[(B)] $1/2$
\item[(C)] $-1/2$
\item[(D)] $-1$
\end{list}

\bigskip\noindent\textbf{A22.}\ $\int_0^1\frac{e^{-x}+1}{e^{-x}}\,dx=$
\begin{list}{}{\setlength{\leftmargin}{3em}\setlength{\labelwidth}{2.5em}\setlength{\itemsep}{0pt}\setlength{\topsep}{2pt}}
\item[(A)] $e$
\item[(B)] $2+e$
\item[(C)] $1/e$
\item[(D)] $1+e$
\end{list}

\bigskip\noindent\textbf{A23.}\ $\int_0^1\frac{e^x}{e^x+1}\,dx=$
\begin{list}{}{\setlength{\leftmargin}{3em}\setlength{\labelwidth}{2.5em}\setlength{\itemsep}{0pt}\setlength{\topsep}{2pt}}
\item[(A)] $\ln2$
\item[(B)] $1+e$
\item[(C)] $-\ln2$
\item[(D)] $\ln\frac{e+1}{2}$
\end{list}

\bigskip\noindent\textbf{A24.}\ If $x=\tan\theta$, then $\int_1^{\sqrt3}\sqrt{1+x^2}\,dx$ is equivalent to
\begin{list}{}{\setlength{\leftmargin}{3em}\setlength{\labelwidth}{2.5em}\setlength{\itemsep}{0pt}\setlength{\topsep}{2pt}}
\item[(A)] $\int_{\pi/4}^{\pi/3}\sec\theta\,d\theta$
\item[(B)] $\int_1^{\sqrt3}\sec^3\theta\,d\theta$
\item[(C)] $\int_{\pi/4}^{\pi/3}\sec^3\theta\,d\theta$
\item[(D)] $\int_{\pi/4}^{\pi/3}\sec^2\theta\tan\theta\,d\theta$
\end{list}

\bigskip\noindent\textbf{A25.}\ If $u=\sqrt{x+1}$, then $\int_0^3\frac{dx}{x\sqrt{x+1}}$ is equivalent to
\begin{list}{}{\setlength{\leftmargin}{3em}\setlength{\labelwidth}{2.5em}\setlength{\itemsep}{0pt}\setlength{\topsep}{2pt}}
\item[(A)] $\int_1^2\frac{du}{u^2-1}$
\item[(B)] $\int_1^2\frac{2\,du}{u^2-1}$
\item[(C)] $2\int_0^3\frac{du}{(u-1)(u+1)}$
\item[(D)] $2\int_1^2\frac{du}{u(u^2-1)}$
\end{list}

\bigskip\noindent\textbf{A26.}\ For the differentiable function $f$, the table is $\begin{array}{crrrrr}x&0&2&4&6&8\\\hline f(x)&11&15&16&12&7\\f'(x)&3&2&-1&-3&0\end{array}$. Evaluate $\int_8^0 f'(x)\,dx$.
\begin{list}{}{\setlength{\leftmargin}{3em}\setlength{\labelwidth}{2.5em}\setlength{\itemsep}{0pt}\setlength{\topsep}{2pt}}
\item[(A)] $-1/2$
\item[(B)] $3$
\item[(C)] $4$
\item[(D)] $-1$
\end{list}

\bigskip\noindent\textbf{A27.}\ Use a midpoint Riemann sum with 3 equal subintervals to approximate the area under $y=6x-x^2$ on $[0,6]$.
\begin{list}{}{\setlength{\leftmargin}{3em}\setlength{\labelwidth}{2.5em}\setlength{\itemsep}{0pt}\setlength{\topsep}{2pt}}
\item[(A)] $19$
\item[(B)] $36$
\item[(C)] $38$
\item[(D)] $54$
\end{list}

\bigskip\noindent\textbf{A28.}\ A continuous function passes through $(4,2),(6,6),(7,5),(10,8)$. Use trapezoids to estimate $\int_4^{10}f(x)\,dx$.
\begin{list}{}{\setlength{\leftmargin}{3em}\setlength{\labelwidth}{2.5em}\setlength{\itemsep}{0pt}\setlength{\topsep}{2pt}}
\item[(A)] $30$
\item[(B)] $32$
\item[(C)] $33$
\item[(D)] $41$
\end{list}

\bigskip\noindent\textbf{A29.}\ Using left and right sums with 2 equal subintervals on $[1,3]$ for $y=1/x$, which gives the resulting lower and upper bounds directly for $\ln3$?
\begin{list}{}{\setlength{\leftmargin}{3em}\setlength{\labelwidth}{2.5em}\setlength{\itemsep}{0pt}\setlength{\topsep}{2pt}}
\item[(A)] $\frac12<\int_1^2\frac{dx}x<1$
\item[(B)] $\frac13<\int_1^3\frac{dx}x<2$
\item[(C)] $\frac13<\int_2^3\frac{dx}x<\frac12$
\item[(D)] $\frac56<\int_1^3\frac{dx}x<\frac32$
\end{list}

\bigskip\noindent\textbf{A30.}\ Let $A=\int_0^1\cos x\,dx$. For left, right, and trapezoidal approximations $L,R,T$ with 100 equal subintervals, which is true?
\begin{list}{}{\setlength{\leftmargin}{3em}\setlength{\labelwidth}{2.5em}\setlength{\itemsep}{0pt}\setlength{\topsep}{2pt}}
\item[(A)] $L<A<T<R$
\item[(B)] $L<T<A<R$
\item[(C)] $R<A<T<L$
\item[(D)] $R<T<A<L$
\end{list}

\bigskip\noindent\textbf{A31.}\ $\int_{-1}^3|x|\,dx=$
\begin{list}{}{\setlength{\leftmargin}{3em}\setlength{\labelwidth}{2.5em}\setlength{\itemsep}{0pt}\setlength{\topsep}{2pt}}
\item[(A)] $7/2$
\item[(B)] $4$
\item[(C)] $9/2$
\item[(D)] $5$
\end{list}

\bigskip\noindent\textbf{A32.}\ $\int_{-3}^2|x+1|\,dx=$
\begin{list}{}{\setlength{\leftmargin}{3em}\setlength{\labelwidth}{2.5em}\setlength{\itemsep}{0pt}\setlength{\topsep}{2pt}}
\item[(A)] $5/2$
\item[(B)] $7/2$
\item[(C)] $11/2$
\item[(D)] $13/2$
\end{list}

\bigskip\noindent\textbf{A33.}\ Find the average value of $y=\sqrt{64-x^2}$ on its domain.
\begin{list}{}{\setlength{\leftmargin}{3em}\setlength{\labelwidth}{2.5em}\setlength{\itemsep}{0pt}\setlength{\topsep}{2pt}}
\item[(A)] $2$
\item[(B)] $4$
\item[(C)] $2\pi$
\item[(D)] $4\pi$
\end{list}

\bigskip\noindent\textbf{A34.}\ Find the average value of $\cos x$ on $[\pi/3,\pi/2]$.
\begin{list}{}{\setlength{\leftmargin}{3em}\setlength{\labelwidth}{2.5em}\setlength{\itemsep}{0pt}\setlength{\topsep}{2pt}}
\item[(A)] $3/\pi$
\item[(B)] $3(2-\sqrt3)/\pi$
\item[(C)] $3/(2\pi)$
\item[(D)] $2/(3\pi)$
\end{list}

\bigskip\noindent\textbf{A35.}\ Find the average value of $\csc^2x$ on $[\pi/6,\pi/4]$.
\begin{list}{}{\setlength{\leftmargin}{3em}\setlength{\labelwidth}{2.5em}\setlength{\itemsep}{0pt}\setlength{\topsep}{2pt}}
\item[(A)] $3\sqrt3/\pi$
\item[(B)] $\sqrt3/\pi$
\item[(C)] $12(\sqrt3-1)/\pi$
\item[(D)] $3(\sqrt3-1)$
\end{list}

\bigskip\noindent\textbf{A36.}\ Choose the Riemann sum limit for $\int_1^5x^3\,dx$.
\begin{list}{}{\setlength{\leftmargin}{3em}\setlength{\labelwidth}{2.5em}\setlength{\itemsep}{0pt}\setlength{\topsep}{2pt}}
\item[(A)] $\lim_{n\to\infty}\sum_{k=1}^n(1+k/n)^3(4/n)$
\item[(B)] $\lim_{n\to\infty}\sum_{k=1}^n(1+4k/n)^3(1/n)$
\item[(C)] $\lim_{n\to\infty}\sum_{k=1}^n(1+4k/n)^3(4/n)$
\item[(D)] $\lim_{n\to\infty}\sum_{k=1}^n(1+k/n)^3(1/n)$
\end{list}

\bigskip\noindent\textbf{A37.}\ Choose the Riemann sum limit for $\int_0^\pi\sin(3x)\,dx$.
\begin{list}{}{\setlength{\leftmargin}{3em}\setlength{\labelwidth}{2.5em}\setlength{\itemsep}{0pt}\setlength{\topsep}{2pt}}
\item[(A)] $\lim_{n\to\infty}\sum_{k=1}^n\sin(3k/n)(1/n)$
\item[(B)] $\lim_{n\to\infty}\sum_{k=1}^n\sin(3\pi k/n)(1/n)$
\item[(C)] $\lim_{n\to\infty}\sum_{k=1}^n\sin(3k/n)(\pi/n)$
\item[(D)] $\lim_{n\to\infty}\sum_{k=1}^n\sin(3\pi k/n)(\pi/n)$
\end{list}

\bigskip\noindent\textbf{A38.}\ Which integral is the limit $\lim_{n\to\infty}\sum_{k=1}^n\sin(1+6k/n)(3/n)$?
\begin{list}{}{\setlength{\leftmargin}{3em}\setlength{\labelwidth}{2.5em}\setlength{\itemsep}{0pt}\setlength{\topsep}{2pt}}
\item[(A)] $\int_1^4\sin(6x)\,dx$
\item[(B)] $\int_0^3\sin(2x+1)\,dx$
\item[(C)] $\int_1^4\sin(2x)\,dx$
\item[(D)] $\int_0^3\sin(6x+1)\,dx$
\end{list}

\bigskip\noindent\textbf{A39.}\ $\int_{\pi/4}^{\pi/2}\sin^3\alpha\cos\alpha\,d\alpha=$
\begin{list}{}{\setlength{\leftmargin}{3em}\setlength{\labelwidth}{2.5em}\setlength{\itemsep}{0pt}\setlength{\topsep}{2pt}}
\item[(A)] $3/16$
\item[(B)] $1/4$
\item[(C)] $-1/4$
\item[(D)] $-3/16$
\end{list}

\bigskip\noindent\textbf{B1.}\ Find the average value of the piecewise linear function joining $(0,0),(2,10),(5,10)$ on $[0,5]$.
\begin{list}{}{\setlength{\leftmargin}{3em}\setlength{\labelwidth}{2.5em}\setlength{\itemsep}{0pt}\setlength{\topsep}{2pt}}
\item[(A)] $4$
\item[(B)] $5$
\item[(C)] $7$
\item[(D)] $8$
\end{list}

\bigskip\noindent\textbf{B2.}\ The integral $\int_{-4}^4\sqrt{16-x^2}\,dx$ gives the area of
\begin{list}{}{\setlength{\leftmargin}{3em}\setlength{\labelwidth}{2.5em}\setlength{\itemsep}{0pt}\setlength{\topsep}{2pt}}
\item[(A)] a circle of radius 4
\item[(B)] a semicircle of radius 4
\item[(C)] a quadrant of a circle of radius 4
\item[(D)] half of an ellipse
\end{list}

\bigskip\noindent\textbf{B3.}\ The graph of $f$ consists of line segments joining $(0,2),(2,4),(6,0)$. Which intervals contain a point $c$ where $f(c)$ equals its average on $[0,6]$? I. $[0,2]$; II. $[2,4]$; III. $[4,6]$.
\begin{list}{}{\setlength{\leftmargin}{3em}\setlength{\labelwidth}{2.5em}\setlength{\itemsep}{0pt}\setlength{\topsep}{2pt}}
\item[(A)] I only
\item[(B)] II only
\item[(C)] III only
\item[(D)] I and II only
\end{list}

\bigskip\noindent\textbf{B4.}\ The graph of $f$ consists of line segments joining $(0,2),(2,4),(6,0)$. Evaluate $\int_0^2f\prime(3x)\,dx$.
\begin{list}{}{\setlength{\leftmargin}{3em}\setlength{\labelwidth}{2.5em}\setlength{\itemsep}{0pt}\setlength{\topsep}{2pt}}
\item[(A)] $-2$
\item[(B)] $-2/3$
\item[(C)] $2/3$
\item[(D)] $2$
\end{list}

\bigskip\noindent\textbf{B5.}\ The graph of $f$ consists of line segments joining $(0,2),(2,4),(6,0)$. Let $g(x)=\int_0^{2x}f(t)\,dt$. Find $g\prime(1)$.
\begin{list}{}{\setlength{\leftmargin}{3em}\setlength{\labelwidth}{2.5em}\setlength{\itemsep}{0pt}\setlength{\topsep}{2pt}}
\item[(A)] $3$
\item[(B)] $4$
\item[(C)] $6$
\item[(D)] $8$
\end{list}

\bigskip\noindent\textbf{B6.}\ The graph of $f$ consists of line segments joining $(0,2),(2,4),(6,0)$. Let $h(x)=x^2-f(x)$. Find $\int_0^6h(x)\,dx$.
\begin{list}{}{\setlength{\leftmargin}{3em}\setlength{\labelwidth}{2.5em}\setlength{\itemsep}{0pt}\setlength{\topsep}{2pt}}
\item[(A)] $38$
\item[(B)] $58$
\item[(C)] $70$
\item[(D)] $74$
\end{list}

\bigskip\noindent\textbf{B7.}\ If $f(x)$ is continuous on the closed interval $[a,b]$, then there exists at least one number $c$, $a<c<b$, such that $\displaystyle\int_{a}^{b} f(x)\,dx$ is equal to
\begin{list}{}{\setlength{\leftmargin}{3em}\setlength{\labelwidth}{2.5em}\setlength{\itemsep}{0pt}\setlength{\topsep}{2pt}}
\item[(A)] $\dfrac{f(c)}{b-a}$
\item[(B)] $f'(c)(b-a)$
\item[(C)] $f(c)(b-a)$
\item[(D)] $\dfrac{f'(c)}{b-a}$
\end{list}

\bigskip\noindent\textbf{B8.}\ If $f(x)$ is continuous on the closed interval $[a,b]$ and $k$ is a constant, then $\displaystyle\int_{a}^{b} kf(x)\,dx$ is equal to
\begin{list}{}{\setlength{\leftmargin}{3em}\setlength{\labelwidth}{2.5em}\setlength{\itemsep}{0pt}\setlength{\topsep}{2pt}}
\item[(A)] $k(b-a)$
\item[(B)] $k[f(b)-f(a)]$
\item[(C)] $\dfrac{[kf(x)]^2}{2}\Big|_a^b$
\item[(D)] $k\displaystyle\int_{a}^{b} f(x)\,dx$
\end{list}

\bigskip\noindent\textbf{B9.}\ $\displaystyle\dfrac{d}{dt}\int_{0}^{t} \sqrt{x^3+1}\,dx =$
\begin{list}{}{\setlength{\leftmargin}{3em}\setlength{\labelwidth}{2.5em}\setlength{\itemsep}{0pt}\setlength{\topsep}{2pt}}
\item[(A)] $\sqrt{t^3+1}$
\item[(B)] $\dfrac{\sqrt{t^3+1}}{3t^2}$
\item[(C)] $3x^2\sqrt{x^3+1}$
\item[(D)] $\sqrt{t^3+1}-1$
\end{list}

\bigskip\noindent\textbf{B10.}\ If $F(u)=\int_1^u(2-x^2)^3\,dx$, find $F\prime(u)$.
\begin{list}{}{\setlength{\leftmargin}{3em}\setlength{\labelwidth}{2.5em}\setlength{\itemsep}{0pt}\setlength{\topsep}{2pt}}
\item[(A)] $-6u(2-u^2)^2$
\item[(B)] $(2-u^2)^3-1$
\item[(C)] $(2-u^2)^3$
\item[(D)] $-2u(2-u^2)^3$
\end{list}

\bigskip\noindent\textbf{B11.}\ Evaluate $\frac{d}{dx}\int_{\pi/2}^{x^2}\sqrt{\sin t}\,dt$.
\begin{list}{}{\setlength{\leftmargin}{3em}\setlength{\labelwidth}{2.5em}\setlength{\itemsep}{0pt}\setlength{\topsep}{2pt}}
\item[(A)] $\sqrt{\sin t^2}$
\item[(B)] $2x\sqrt{\sin x^2}-1$
\item[(C)] $\sqrt{\sin x^2}-1$
\item[(D)] $2x\sqrt{\sin x^2}$
\end{list}

\bigskip\noindent\textbf{B12.}\ On the upper ellipse branch $x=4\cos\theta$, $y=3\sin\theta$, $0\leq\theta\leq\pi/3$, the integral $\int_2^4xy\,dx$ is equivalent to
\begin{list}{}{\setlength{\leftmargin}{3em}\setlength{\labelwidth}{2.5em}\setlength{\itemsep}{0pt}\setlength{\topsep}{2pt}}
\item[(A)] $48\int_{\pi/3}^0\sin\theta\cos^2\theta\,d\theta$
\item[(B)] $48\int_2^4\sin^2\theta\cos\theta\,d\theta$
\item[(C)] $48\int_0^{\pi/3}\sin\theta\cos^2\theta\,d\theta$
\item[(D)] $48\int_0^{\pi/3}\sin^2\theta\cos\theta\,d\theta$
\end{list}

\bigskip\noindent\textbf{B13.}\ Use a left rectangular sum on the five table intervals to estimate $\int_0^{30}f(x)\,dx$: $\begin{array}{c|rrrrrr}x&0&5&12&20&24&30\\\hline f(x)&9&11&7&3&1&2\end{array}$.
\begin{list}{}{\setlength{\leftmargin}{3em}\setlength{\labelwidth}{2.5em}\setlength{\itemsep}{0pt}\setlength{\topsep}{2pt}}
\item[(A)] $144$
\item[(B)] $170$
\item[(C)] $186$
\item[(D)] $196$
\end{list}

\bigskip\noindent\textbf{B14.}\ At which $x\in[0,10]$ does $y=x^2$ reach its average value on that interval?
\begin{list}{}{\setlength{\leftmargin}{3em}\setlength{\labelwidth}{2.5em}\setlength{\itemsep}{0pt}\setlength{\topsep}{2pt}}
\item[(A)] $5.313$
\item[(B)] $5.774$
\item[(C)] $6.083$
\item[(D)] $18.257$
\end{list}

\bigskip\noindent\textbf{B15.}\ Find the average of $f(x)=\begin{cases}x^3&x<2\\4x&x\geq2\end{cases}$ on $[0,5]$.
\begin{list}{}{\setlength{\leftmargin}{3em}\setlength{\labelwidth}{2.5em}\setlength{\itemsep}{0pt}\setlength{\topsep}{2pt}}
\item[(A)] $6.25$
\item[(B)] $6.41$
\item[(C)] $9.2$
\item[(D)] $10$
\end{list}

\end{document}
